%\documentclass[14pt,tlenc]{article}

\documentclass[12pt]{article}
%\Large
\include{usepack}
%\include{my_environ}
%for version counter



% environment for computer listing

\renewcommand{\thefootnote}{\fnsymbol{footnote}}


%commands for shortening the typing

\newcommand{\polp}{$p^{\uparrow}$}

\newcommand{\pold}{$d^{\uparrow}$}

\newcommand{\xT}{$x_{2}$}

\newcommand{\xB}{$x_1$}

\newcommand{\xF}{$x_F$}

\newcommand {\poc}{POC}

\newcommand{\ra}{$\rightarrow$}

%\newcommand{\bl}{\textcolor{blue}}
%\newcommand{\bl}[1]{\textcolor{blue}{#1}}

% define counter for lab
\newcounter{labcount}
%\makeindex
\begin{document}



\author{ Andi Klein}
\title{ MyCal \\ 
		\bf{version: 0.5}}
%\date{12 Nov 2002}
%\today
\maketitle
% set the counter level depth
\setcounter{secnumdepth}{10}
\setcounter{tocdepth}{10}



\newenvironment{andilist}{\begin{itemize} \em}{\end{itemize}}
% end of definitions
\begin{abstract}
MyCal is a program to calculate calories from a database of ingredients. It is based on {\bf psql} and {\bf pyside6}
\end{abstract}
\pagebreak

\tableofcontents
%\section{Tasks to be done}
%\begin{enumerate}
%\item check references. Make sure xarch is replaced with real pub
%\end{enumerate}
%\pagebreak

\section{Introduction}
In the following I will try to descripbe the setup of the program, so that in a few month I still can remember what I did.

\section{Layout of the program}
Currently, the main class is ControlDB, which has the methods described below. While in the beginning I created the form for ingredients from scratch, I later started to use the QT designer. The advantage of the QT designer is that I can easily add or move things in a design, run the conversion from ui to py and be done.
Calling the designer is done by {\bf pyside6-designer Recipe.ui \&} to open the Recipe.ui form. After having changed the design and saved it, running {\bf pyside6-uic Recipe.ui -o Recipe.py} converts the design into a python class, which can the be reused.


\section{class NewWindow}
Creates a simple mainwindow, which can then be used to add widgets.

\section{class RecipeModel}
This is created from QAbstractListModel and creates the model for the recipe form. This is part of the model view architecture of pyside6 and is used by the method {\bf MyRecipes\_new}.

\section{class ControlDB}
This the main class of the program. It is called by
\begin{verbatim}

class ContrlDB(QMainWindow):

    def __init__(self,Title=None,db_name=None,db_user=None,db_system=None,db_pwd = None):8
        super().__init__()
        self.db_name = db_name
        self.db_user = db_user
        self.db_system = db_system
        if db_pwd == None:
\end{verbatim}

where the database parameters are passed.

\subsection{Important variables}
\begin{enumerate}
\item self.panda\_ingredients[][] list of ingredients in a recipe
\item self.MRD.recipes The table entry of the recipes. This contains the list of weights and ingredients
\item self.selected\_recipe name of the recipe clicked on 



\end{enumerate}

\subsection{some poogrammatic issues}
The connection to the tableviewmodel goes through {\bf self.RecipeModel.recipes = []}. Defining this creates the data in the model. 
\subsection{ConnectDataBase}
This establishes the connection to the psql database, using a text file for the password. 
If the connection fails , the program exits

\subsection{SetupLogger}
This creates the logger system. It deals with warnings and errors.

\subsection{ShowTables}
Lists all the tables in the database

\subsection{ViewTable}
This method sets up the connection to a specific table and you can view this table. You can choose the list of columns which you want to suppress in display.
\begin{verbatim}
def ViewTable(self,table,suppress_columns=[],editmode=True, 
columnwidth=[],Title = None, showTable = True):
\end{verbatim}
Calling this method with showTable = False will suppress the display. This is one of the main methods to deal with tables and selecting them.

\subsection{SetPosition}
Move the window to position chosen.

\subsection{CreateIngredientsForm}
This creates the form to add ingredients. This builds the respective form from scratch, something I don't want to do .
It connects to SaveIngredients and also to Cancel

\subsection{InsertRecord}
First it creates a tableview with the above ViewTable for the ingredients table and then creates a QSqlQuerymodel. It inserts self.record, which has been created by SaveIngredients, which in turn is a connected method called by clicking on the SaveButton from the ingredient form.

\subsection{MyRecipes\_new}
This is the heart of the communication with the recipes table. This establishes the recipe window creates the RecipeView, as well as the RecipeModel. It calls the GetRecipeList, which creates a list of all the recipe names.

\subsection{getSelectedText, getSelectedRecipes}
This is a function called through a connector in the Recipe Combobox. First it clears the record, and then 
fills the record with the table row chosen through the combo box. It uses an emit signal to refresh the list.
\subsection{FillRecipeIngredientList}
This is currently fairly long for historical reasons. The columns , which contains the ingredients of a specific recipe from the original Libreoffice days looks like this:

\begin{verbatim}
177 g Flour unbleached *u000a79 g Butter *u000a33 g Almond flour *u000a480 g Apricots frozen  *u000a8 g Sugar *
\end{verbatim}
There are a few problems with this:
\begin{enumerate}
\item We have spaces in the name of some of the ingredients. I manually changed the ingredients table to have always underscores, but the recipe table is too long.
\item We have characters in which have to do with textformatting (*u000a)
\item It is one string
\end{enumerate}

In this routine we strip out the UTF characters, remove the *, and then create a list in self.recipe\_ingredients

{\bf Note in vs03 more handling} create 2 dimensional list because now we have tableview
\subsection{DisplayRecipeList}
Adds items to the Recipelist.

\subsection{split\_and\_keep}
This is used by FillRecipeIngredientList to separate the long string into a list of values, where the separation character is {\bf g}.
\subsection{StripUnicode}
removes the u000a from the record.

\subsection{GetRecipeList}
Gets a list of all the recipes in the recipe table.

\subsection{SizeWindow}
Creates dimensions of the window selected. it gives the upper corner coordinate and the x and y size.

\subsection{SaveIngredients}
Connected to from the CreateIngredientsForm, it checks that the new ingredient is unique and then saves the new created ingredient in the table.

\subsection{on\_item\_click}
gets connected to QlistView, and called when you click on one of the items.

\subsection{save\_recipe}
Called from save button in recipe form, it saves the current record for ingredients. It then calls pack\_record, so that we can update the record.

\subsection{pack\_record}
The record from save\_recipe needs to be packed back into a single string.

\subsection{update\_recipe}
called from save recipe it updates the ingredient record.

\subsection{update\_record}
This updates   row b with name in table

\subsection{do\_sql}
Executes a sql statement whihc has been formed by the calling method.

\subsection{create\_pandas\_frame}
This packs all the pertinent information for CalculateCalories into a pandas frame.


\section{Some Useful commands}
\begin{enumerate}
\item {\bf pyside6-designer Recipe.ui \&}
\item {\bf pyside6-uic Recipe.ui -o Recipe.py}
\end{enumerate}

\section{Notes}
If you change a value in the table form, you have to hit a <CR>; same for entering the name.

It is important to only use underscores in names or ingredients. Do not use spaces. All entries are in {\bf gram}. I find the US system of oz, cups etc totally stone age; so will never use that system.

do {\bf NOT} use plural for any quantity

\section{Current Status}
I think with QListview for the recipes, I boxed myself into a corner. This is the end of vs 0.2
I create now vs 0.3 where the interface is a tableview. Tableview is working, next step is to be able to edit the fields and then save them.

\section{Current issues}
\begin{itemize}
\item The unpacking of a ingredienst list with the new stored value only gives the first entry. Needs to be solved first. fixed
\item finish new recipe setup, current problem is it does not strip the empty records at the end.
\item menu
\item Populate the textfield in the form like cal\/gram etc
\item Create the calorie calculator
\item Make sure entries in table get checked for upper and lower characters
\item Test for existence in ingredients database. Handle the case where it does not exist.
\item Possibly use US govt nutrition database. https://fdc.nal.usda.gov/
\item Clean up the code with consistent name scheme.
\item ingredient names plural and singular. set all to singular
\item Currently it looks like the UPDATE actually deletes the record. This might have to do with my do\_sql changes. 
Need backup of psql table. actually it does not delete it but also not update it. again this might be my do\_sql.

\item Creating new recipe takes the old food database, editing an existing one takes the new. At least in the menu version. Check if this is also the case with controldb\_new. yes it is in controldb.
\item still crashes line 871 with duplicate entry, also crashes with 769 with new\_recipe\_record does not exist

\end{itemize}

\section{Design of Calculation}
Before I start, I want to think about the design of the center piece, calculating the calories etc of the recipes. In order to do this, I will need to create a new class. The program as it stands in vs .3 is already unwieldy enough and testing is pretty darn difficult. So this should be a stand alone class.

In the following I will sketch some of the stuff needed.
\begin{enumerate}
\item  For every recipe I need to access both databases. To get the recipes, ingredients and amounts as well as 
the values for each ingredient.
\item I will use pandas to deal with the math and storing the variables. Manipulating the variables is more straight forward.
\item I will try to keep this class completely free from pyside6, so that I can keep the code as clean as possible. This will require some kind of interface between the two. This will provide the data from ControlDB to the calculate class. I think the best way is to already fill a data frame in ControlDB and then pass this dataframe to the CalculateCalories class.
\item In the controlDB, I still need to create a function, which will deal with a missing ingredient. Shall I automatically open the fill ingredient form? 
\item Since potentially I would like to use the US nutrition database, I might be able to access a lot more info than just the calories. So I will try to design the class accordingly.
\item In order 


\end{enumerate}
\end{document}
